% =========================================================
% PACKAGE DECLARATIONS
% =========================================================
% All LaTeX packages required by this template are loaded here.
% Each package includes a comment explaining why it is needed.
%
% IMPORTANT: Do not remove packages unless you are sure they
% are not used anywhere in the template. Adding new packages
% here is safe — just place them at the end of this file.
% =========================================================

% ---------------------------------------------------------
% FONT SUPPORT (XeLaTeX)
% ---------------------------------------------------------
% fontspec: Required for XeLaTeX to use system fonts like
% Times New Roman. This is the recommended approach for
% matching the Anna University requirement exactly.
% DO NOT REMOVE — the entire template depends on this.
\usepackage{fontspec}

% ---------------------------------------------------------
% LANGUAGE AND ENCODING
% ---------------------------------------------------------
% For English-only documents with XeLaTeX, explicit language
% setup is not required. If you need multilingual support,
% uncomment the two lines below:
% \usepackage{polyglossia}
% \setmainlanguage{english}

% ---------------------------------------------------------
% PAGE LAYOUT
% ---------------------------------------------------------
% geometry: Controls page dimensions and margins.
% The margin values here implement the official Anna University
% requirements. See FORMAT_RULES.md, Section 2.
\usepackage{geometry}

% ---------------------------------------------------------
% LINE SPACING
% ---------------------------------------------------------
% setspace: Provides \onehalfspacing, \doublespacing,
% and \singlespacing commands for controlling line spacing
% throughout the document.
\usepackage{setspace}

% ---------------------------------------------------------
% HEADERS AND FOOTERS
% ---------------------------------------------------------
% fancyhdr: Customizes headers and footers.
% Used to place page numbers at the center bottom.
\usepackage{fancyhdr}

% ---------------------------------------------------------
% TABLE OF CONTENTS FORMATTING
% ---------------------------------------------------------
% tocloft: Customizes the Table of Contents, List of Figures,
% and List of Tables formatting (fonts, spacing, leader dots).
\usepackage[titles]{tocloft}

% ---------------------------------------------------------
% CHAPTER AND SECTION HEADING STYLES
% ---------------------------------------------------------
% titlesec: Customizes the formatting of chapter titles,
% section headings, and subsection headings to match the
% Anna University format (centered chapter titles, etc.).
\usepackage{titlesec}

% ---------------------------------------------------------
% GRAPHICS
% ---------------------------------------------------------
% graphicx: Required for inserting images (\includegraphics).
% Used for logos, diagrams, screenshots, and figures.
\usepackage{graphicx}

% ---------------------------------------------------------
% COLORS (optional, for hyperlinks)
% ---------------------------------------------------------
\usepackage{xcolor}

% ---------------------------------------------------------
% TABLES
% ---------------------------------------------------------
% booktabs: Provides professional-quality table rules
% (\toprule, \midrule, \bottomrule) instead of ugly \hline.
\usepackage{booktabs}

% array: Extended column definitions for tables.
\usepackage{array}

% longtable: For tables that span multiple pages.
\usepackage{longtable}

% tabularx: Tables with auto-adjusting column widths.
\usepackage{tabularx}

% multirow: Allows cells to span multiple rows in tables.
\usepackage{multirow}

% ---------------------------------------------------------
% CAPTIONS
% ---------------------------------------------------------
% caption: Customizes figure and table caption formatting.
% Used to make captions bold and properly sized.
\usepackage{caption}

% ---------------------------------------------------------
% MATHEMATICAL TYPESETTING
% ---------------------------------------------------------
% amsmath: Standard package for mathematical equations,
% aligned environments, and math symbols.
\usepackage{amsmath}

% amssymb: Additional math symbols.
\usepackage{amssymb}

% ---------------------------------------------------------
% CODE LISTINGS
% ---------------------------------------------------------
% listings: For typesetting source code with syntax highlighting.
% Used in Chapter 6 (Results and Coding) for sample code.
\usepackage{listings}

% ---------------------------------------------------------
% HYPERLINKS AND CROSS-REFERENCES
% ---------------------------------------------------------
% hyperref: Makes cross-references, citations, TOC entries,
% and URLs clickable in the PDF. Colors are set to black
% for print compatibility.
\usepackage[hidelinks]{hyperref}

% ---------------------------------------------------------
% BIBLIOGRAPHY
% ---------------------------------------------------------
% natbib: Flexible bibliography support with author-year
% and numbered citation styles. Used with plainnat style
% for alphabetical ordering as required by AU guidelines.
\usepackage[numbers,sort&compress]{natbib}

% ---------------------------------------------------------
% APPENDIX SUPPORT
% ---------------------------------------------------------
% appendix: Provides additional control over appendix formatting.
\usepackage[toc,page]{appendix}

% ---------------------------------------------------------
% LIST CUSTOMIZATION
% ---------------------------------------------------------
% enumitem: Customizes itemize, enumerate, and description lists.
\usepackage{enumitem}

% ---------------------------------------------------------
% PDF METADATA
% ---------------------------------------------------------
% This sets the PDF document properties (visible in PDF readers).
% The values are pulled from metadata.tex commands.
\AtBeginDocument{%
  \hypersetup{
    pdftitle={\ProjectTitle},
    pdfauthor={\StudentOneName},
    pdfsubject={B.E. Project Report},
    pdfkeywords={Anna University, Project Report, \BranchName},
  }%
}
